\section{Resumen}

Este trabajo analiza la viabilidad de dotar a los asistentes virtuales de un nivel de hiperpersonalización que combine datos semánticos y estructurados del usuario. Plantea la hipótesis de que, a partir del análisis de interacciones previas y la ingestión de documentos relevantes, es posible construir una base de datos híbrida capaz de aportar contexto en tiempo real a la generación de respuestas.

Para contrastar esta premisa se diseñó una metodología experimental en cuatro etapas:

\begin{enumerate}
\item Experimentos de \textit{prompting} acotados
\item Ingestión de documentos para ampliar el corpus contextual
\item Generación de conjuntos de datos sintéticos
\item Almacenamiento de dichos conjuntos en bases vectoriales y de grafos, con filtros específicos por usuario
\end{enumerate}

El prototipo, que integra Redis para sesiones y FAISS para búsquedas por similitud, demostró que la distancia coseno, aplicada tras un filtrado por identificador, recupera información pertinente aun en corpora crecientes; el sistema resolvió consultas cronológicas y semánticas combinadas con alta precisión.

Se concluye que la estrategia es técnicamente factible y escalable, siempre que el motor vectorial gestione eficientemente el volumen de datos; no obstante, persisten retos ligados a la memoria conversacional continua y a la adaptación de modelos menos potentes.